\begin{topic}[id=embedded_vision_ondevice,
             difficulty={Mittel},
             type={Systemanalyse + Benchmarking},
             prerequisites={Grundlagen eingebetteter Systeme; Basiswissen Bildverarbeitung oder Python/C++},
             practice={gut möglich},
             owner={Dozententeam},
             updated={2026-04-09},
             tags={ Embedded Vision, OpenCV, Jetson, Pipeline, FPS, Latenz, Preprocessing, Optimierung }]{Embedded Vision on-device (z.B. Jetson/MCU)}

\TopicDescription{On-device-Vision reduziert Latenz und schützt Daten, erfordert aber sorgsame Pipeline-Gestaltung: Kamera-I/O, Farbraum-Konvertierung, Resizing und Pufferverwaltung sind oft die wahren Bottlenecks. Das Thema zeigt robuste Minimal-Pipelines auf Edge-Geräten und kombiniert klassische Verfahren mit leichten Netzen. Du lernst Metriken (FPS, End-to-End-Latenz, Auslastung) sinnvoll zu nutzen und einfache Optimierungen (ROI-Cropping, Quantisierung, Preprocessing am Sensor) systematisch zu bewerten.}

\TopicLeitfragen{\item Wo entstehen in Edge-Vision-Pipelines die größten Latenzen?
\item Welche einfachen Optimierungen bringen messbaren Nutzen?
\item Wann lohnt sich ein Hardware-Beschleuniger?
}
\TopicScope{\item Kein Training großer Netze.
\item Keine umfangreiche Kameratreiber-Entwicklung.
}
\TopicPractice{Mini-Pipeline „Kamera → Preprocessing → leichte Inferenz/Heuristik → Ausgabe“ mit zwei Optimierungsvarianten.}

\TopicKeywords{Embedded Vision, OpenCV, Jetson, Pipeline, FPS, Latenz, Preprocessing, Optimierung}

\nocite{opencvDocs,nvidiaJetsonGuide}
\end{topic}
